\section{Proofs and additional results}
\begin{frame}[t,label=singleton]{1. A simpler gap with one client per firm}
\small
Start each of $N$ independent copies with
\[
 H=\{ax,ay,az,bx,by,bz,cz\}.
\]
$\{ay,bx,cz\}$ is already supported and full-market stable.
DA instead returns $\{ax,by,cz\}$, blocked by $cx$.

\medskip
Without auditing $c$ or $x$, the only possible addition is $cy$.
It changes nothing, since $c$ keeps $z$. Auditing either $c$ or $x$ repairs the copy.

\medskip
Add one hidden isolated pair. With unit costs and private firms,
\[
 \supp=1,\qquad \resp=N+1.
\]
\sources{The shared-firm construction also shows how audits needed in one part of
the market can destabilize another part.}
\end{frame}

\begin{frame}[t,label=master]{2. The support master in equations}
\fontsize{11}{13}\selectfont
$x_f\in\{0,1\}$ chooses audits. Exposure $v_{bs}\in[0,1]$ obeys
\[
 v_e=1\ (e\in E_0),\qquad
 v_{bs}\geq x_{f(b)},x_{f(s)},\quad
 v_{bs}\leq x_{f(b)}+x_{f(s)}\ (bs\notin E_0).
\]
Let $\delta(i)$ denote incident pairs. Minimize $\sum_fc_fx_f$ with $z_e\geq0$:
\begin{align*}
 \sum_{e\in\delta(i)}z_e&\leq1 &&\text{for every client }i,\\
 z_{bs}+\sum_{s'\succ_b s}z_{bs'}+\sum_{b'\succ_s b}z_{b's}&\geq1 &&bs\in E,\\
 z_e&\leq v_e &&e\in E.
\end{align*}
For fixed integral $x$, $v$ is integral. Setting hidden $z_e$ to zero gives a face
of the integral stable-matching polytope, so continuous $z$ is enough.
\sources{Rothblum (1992). This argument is for a fixed integral audit vector.
It does not assert integrality when the audit variables are fractional.}
\end{frame}

\begin{frame}[t,label=replay]{3. Why the prefix constraint is valid}
\small
Replay DA in the same sequential proposer order. If all full-prefix exposure
statuses agree, the next buyer finds the same available choice. The seller
makes the same comparison. Induction preserves every tentative match.

\medskip
For a failed component $C$ with edge set $E_C$, let
$P=(\Pi(M)\cap E_C)\setminus E_0$. If the current exposure is $v^t$, add
\[
 \sum_{e\in P:v_e^t=0}v_e+
 \sum_{e\in P:v_e^t=1}(1-v_e)\geq1.
\]
Every successful plan must change one of these statuses, or the same local
outcome and its blocking pair persist.
\sources{An unmatched buyer contributes its whole acceptable list. Full-market components
are independent for matching, even when firms couple their audit decisions.}
\end{frame}

\begin{frame}[t,label=exactness]{4. Exactness and stopping}
\small
Let $\FF_C$ be the firms with clients in component $C$. Every successful plan
stays feasible for the master. Each failure excludes a new audit assignment.
\[
 \begin{array}{ll}
 \text{Whole-assignment method:}&2^{|\FF|}+1\text{ master solves at most},\\[3pt]
 \text{Local methods:}&1+\displaystyle\sum_C2^{|\FF_C|}\text{ at most}.
 \end{array}
\]
A verified master candidate is optimal because its cost equals a valid
lower bound. The same is true when that bound reaches a verified incumbent.

\medskip
On interruption, use the master's certified dual bound as $L$.
A candidate supplies $U$ only after DA passes the full-market check.
\sources{The theorem assumes exact master optimization and exact oracle evaluation.
Implemented MILP bounds have numerical tolerances. The bounds can be exponential.}
\end{frame}

\begin{frame}[t,label=round-proof]{5. Two solves versus exponentially many failures}
\small
The shared-firm instance has a unique support optimum $\{F_A,F_B\}$.
Its first failed response produces $N$ constraints
\[
 (1-v_{a_ix_i})+(1-v_{b_iy_i})+v_{c_ix_i}\geq1.
\]
Any competitive plan contains $F_A,F_B$, so these reduce to $v_{c_ix_i}=1$.
The second master therefore pays one private repair per copy and certifies $N+2$.

\medskip
Every set $\{F_A,F_B\}\cup\{f(c_i):i\in T\}$ with $T\subsetneq[N]$ supports stability,
but DA fails and its cost is below $N+2$. There are $2^N-1$ such sets.
Whole-assignment cuts must remove each separately.
\sources{Both methods start with the support master and a verified upper bound.
The comparison excludes initialization rounds and applies only to the specified cut families.}
\end{frame}

\begin{frame}[t,label=prefix-cover]{6. A sufficient certificate can cost too much}
\small
Let $M^*=\DA(E)$. Exposing every edge in $\Pi(M^*)\setminus E_0$
reproduces $M^*$ by replay. Paying for these edges is vertex cover on endpoint firms.

\medskip
But in this two-by-two component,
\[
 a:s\succ r,\quad b:s\succ r,\quad
 s:a\succ b,\quad r:b\succ a,
\]
$E_0=\{as,br\}$ already gives $M^*$, while the full prefix contains hidden $bs$.
That extra disclosure is unnecessary for the outcome.

\medskip
Take $N$ private copies and one separate switch needing one audit.
The true optimum is one. Prefix cover costs $N+1$.
\sources{This also separates prefix completeness from reproducing $M^*$ itself.
It limits the certificate's cost, not its correctness.}
\end{frame}

\begin{frame}[t,label=tractable]{7. Two other tractable cases}
\small
\textbf{Switch markets with small portfolios.}
With private sellers, unit costs, distinct endpoint firms, and no repeated
firm pair, two clients per buying firm give a degree-two firm graph.
Vertex cover on its paths and cycles is linear-time.

\medskip
\textbf{Bounded interaction width.}
For each full-market component $C$, tabulate the local success predicate
$\phi_C(x|_{\FF_C})$. Make $\FF_C$ a clique in a graph of firms.
Given a width-$t$ tree decomposition, Boolean constraint DP takes
\[
 O\!\left(\sum_C2^{|\FF_C|}\operatorname{poly}(|C|)
       +2^{t+1}\operatorname{poly}(|I|)\right).
\]
\sources{The first result applies only to the switch subclass. The second uses standard
constraint dynamic programming. Here $|I|$ is input length. A component containing many firms can make $t$ large.}
\end{frame}

\begin{frame}[t,label=pitfalls]{8. Why simple repair rules can fail}
\small
\textbf{Auditing a blocker's endpoints is not necessary.}
In the switch, $qy$ initially blocks. Auditing $p$'s firm removes the block
by producing $py,qx$, although that firm represents neither $q$ nor $y$.

\medskip
\textbf{Blocker reduction need not be submodular.}
Let $a,b$ both rank $x\succ y$, and $x,y$ both rank $a\succ b$.
Start with $E_0=\{bx\}$ and private firms.
\begin{center}
\begin{tabular}{@{}lrrrr@{}}
\toprule
Audited buyers & $\varnothing$ & $a$ & $b$ & $a,b$\\
\midrule
Number of blockers & 2 & 1 & 2 & 0\\
\bottomrule
\end{tabular}
\end{center}
For reduction $g(A)$, $g(a)+g(b)=1<2=g(a,b)+g(\varnothing)$.
\end{frame}

\begin{frame}[t,label=bonus]{9. A separate model for partnership bonuses}
\small
All pairs in $E$ are available, but recommendation scores change:
\[
 \widetilde u_i(e)=u_i(e)+\lambda_iF_e,\qquad
 F_e\in\{0,1\},\quad\lambda_i\geq0.
\]
Both endpoints share the same partnership label $F_e$.
Scores and bonuses are rational. Let $P,Q$ be the strict base and modified
rankings, with the outside option last.

\medskip
Define $x\succ_i^*y$ when $i$ prefers $x$ to $y$ in \emph{both} rankings.
Then
\[
 \text{stable under both }P,Q
 \quad\Longleftrightarrow\quad
 \text{super-stable under }\succ^*.
\]
Existence and construction are polynomial-time by partial-order super-stability.
\sources{This is a separate model from the access problem. The shared binary label
enables the reduction to Cseh and Juhos (2021).}
\end{frame}

\begin{frame}[t,label=bonus-proof]{10. Why the binary label matters}
\small
For a candidate pair $e$ and endpoint $i$, let $e_i^M$ be its current match.
Give the unmatched option label zero. The gain from deviating satisfies
\[
 \Delta_i^Q(e)=\Delta_i^P(e)+\lambda_i(F_e-F_{e_i^M}).
\]
If $F_e=1$, both endpoints' gains weakly increase. Stability in $Q$ supplies
an endpoint that rejects $e$ in $Q$, and therefore in $P$ too.

\medskip
If $F_e=0$, both gains weakly decrease. Use the rejecting endpoint from $P$.
Thus every unmatched pair has one endpoint rejecting it under both profiles,
which is exactly super-stability of the intersection orders.
\sources{Conversely, a common rejecting endpoint prevents blocking in each profile.
The same rejection holds for independent strengths $\theta_i\in[0,\lambda_i]$
because each score difference is affine.}
\end{frame}

\begin{frame}[t,label=bonus-limit]{11. The bonus result has limits}
\small
A simultaneous solution need not exist. In a two-by-two market, give everyone
a base score of one for their diagonal partner and zero for the other.
Give crossed pairs a bonus of two. The two profiles have different unique stable matchings.

\medskip
Three label levels also break the binary reduction. Take
$M=\{ar,ts\}$ and candidate $as$, with
\[
 F_{ar}=2,\quad F_{as}=1,\quad F_{ts}=0,\qquad\lambda_i=2.
\]
Set $u_a(r)=0,u_a(s)=1$ and $u_s(a)=0,u_s(t)=1$.
Seller $s$ rejects $as$ under base scores. Buyer $a$ rejects it after bonuses.
There is no common rejecting endpoint.
\sources{This disproves that reduction for three levels. It does not prove NP-hardness.
Use the three-edge acceptable graph $\{ar,as,ts\}$, with sole partners for $r,t$.}
\end{frame}

\begin{frame}[t,label=repro]{12. What was recorded and checked}
\small
\begin{itemize}
\setlength{\itemsep}{5pt}
\item 160 core, 32 larger, and 32 cost/portfolio markets.
The 40 ablation runs reuse core instances.
\item Eight constructed response-gap sizes and five prefix-cover-gap sizes,
kept separate from random-market results.
\item 1,280 core method outputs and 614 optimal plans independently replayed.
DA, prefix replay, reductions, and small optima checked against enumeration.
\item CPU runs on ARM64 macOS, one HiGHS thread, fixed seeds,
rotated method order, source hashes, and explicit time caps.
\end{itemize}
\sources{95\% intervals use 20,000 bootstrap replicates over 20 shared seed blocks.
All final core matchings are perfect because the generator uses balanced complete blocks.
That is a generator property, not evidence about housing scarcity.}
\end{frame}

\begin{frame}[t,label=priority]{13. A common priority removes the gap}
\small
Suppose all sellers use the same strict priority over buyers.
There is a unique stable matching $M^*$: clients choose in priority order.
\[
 M^*\subseteq H\subseteq E\quad\Longrightarrow\quad\DA(H)=M^*.
\]
In any visible graph $H$ containing $M^*$, each full-market choice remains
available at its turn. Hence
\[
 \supp=\resp.
\]
We just need to cover the hidden edges of $M^*$ with endpoint-firm audits.
For side-exclusive firms, weighted bipartite vertex cover solves this by minimum cut.
\sources{A consequence of serial dictatorship. The minimum-cut step requires side-exclusive firms.
Other tractable cases appear in Appendix 7.}
\end{frame}

\begin{frame}[t,label=study]{14. Proposed listing study}
\small
\begin{columns}[T,onlytextwidth]
\begin{column}{.48\textwidth}
\textbf{What we could observe}\par
Rents and features in a fixed listing catalogue.\par
Listing IDs recommended by a fixed AI model.\par
Census household distributions for generating fictional profiles.
\end{column}
\begin{column}{.48\textwidth}
\textbf{What would remain modeled}\par
Preferences on both sides.\par
Agent portfolios and audit costs, unless independently verified.\par
The matching rule and its resulting allocations.
\end{column}
\end{columns}
\medskip
Hold catalogue and recommendation count fixed. Randomize listing order,
repeat prompts, and compare audit policies on the same markets.
\sources{Proposed, not run. Candidate sources: Two Sigma rental listings and American Community Survey microdata.
Verify dates and manager-ID meanings. Inquiry interest is not a preference list.
Do not link Census records to individual listings.}
\end{frame}

\begin{frame}[t,label=access]{15. The cost of an access target}
\small
Fix each modeled buyer's affordable feasible options $F_b$.
Let $G^+$ be lower-income households with at least one such option.
\[
 \mathrm{Access}(A)=\frac{1}{|G^+|}\sum_{b\in G^+}
 \frac{|\{s\in F_b:bs\in E(A)\}|}{|F_b|}.
\]
Let $C(\tau)$ be the cheapest plan that makes DA full-market stable
and achieves $\mathrm{Access}(A)\geq\tau$. Study
\[
 \Delta C(\tau)=C(\tau)-C(0).
\]
How does that extra cost change with scarcity, firm concentration,
and the recommender that chooses $E_0$?
\sources{A proposed extension, not a completed result. Assume $|G^+|>0$ and
$\tau\in[0,1]$. Report households with no affordable option separately.
Access measures options, not eventual leases.}
\end{frame}

\begin{frame}[t,label=refs-theory]{16. Matching and optimization references}
\fontsize{9.5}{11.5}\selectfont
\begin{enumerate}
\setlength{\itemsep}{3pt}
\item \href{https://doi.org/10.1080/00029890.1962.11989827}{Gale and Shapley (1962), College Admissions and the Stability of Marriage}.
\item \href{https://doi.org/10.1016/j.disopt.2016.03.002}{Cseh and Manlove (2016), Stable Marriage and Roommates Problems with Restricted Edges}.
\item \href{https://arxiv.org/abs/1602.04792}{Rastegari, Goldberg, Manlove (2016), Preference Elicitation in Matching Markets via Interviews: A Study of Offline Benchmarks}.
\item \href{https://doi.org/10.1613/jair.1.12755}{Boehmer et al.\ (2021), Bribery and Control in Stable Marriage}.
\item \href{https://arxiv.org/abs/2609.09693}{Aguiar and Hong (2026), The Attention Cost of Stable Matching}. Preprint.
\item \href{https://doi.org/10.1007/s10107-003-0375-9}{Hooker and Ottosson (2003), Logic-Based Benders Decomposition}.
\href{https://doi.org/10.1007/BF01586041}{Rothblum (1992), Characterization of Stable Matchings as Extreme Points of a Polytope}.
\item \href{https://doi.org/10.1016/0304-3975(76)90059-1}{Garey, Johnson, Stockmeyer (1976), Some Simplified NP-Complete Graph Problems}.
\href{https://doi.org/10.1016/S0304-3975(98)00158-3}{Alimonti and Kann (2000), Some APX-Completeness Results for Cubic Graphs}.
\item \href{https://hpi.de/friedrich/docs/publications/2021/TEAC.pdf}{Cseh and Juhos (2021), Pairwise Preferences in the Stable Marriage Problem}.
\end{enumerate}
\end{frame}

\begin{frame}[t,label=refs-housing]{17. Housing, platforms, and data}
\fontsize{10.5}{13}\selectfont
\begin{enumerate}
\setlength{\itemsep}{5pt}
\item \href{https://arxiv.org/abs/2508.18325}{Trabelsi et al.\ (2025), Facilitating Matches on Allocation Platforms}. ECAI.
\item \href{https://r-t-lab.org/works/housing-gpt-audit/}{Liu, So, Hosoi, D'Ignazio (2024), Racial Steering by Large Language Models}. EAAMO.
\item \href{https://doi.org/10.1257/app.20160214}{Barwick, Pathak, Wong (2017), Conflicts of Interest and Steering in Residential Brokerage}. AEJ: Applied Economics.
\item \href{https://ojs.aaai.org/index.php/ICWSM/article/view/7276}{Asplund et al.\ (2020), Auditing Race and Gender Discrimination in Online Housing Markets}. ICWSM.
\item \href{https://doi.org/10.1215/00703370-9357518}{Besbris, Schachter, Kuk (2021), The Unequal Availability of Rental Housing Information Across Neighborhoods}. Demography.
\item Proposed data: \href{https://www.kaggle.com/competitions/two-sigma-connect-rental-listing-inquiries/data}{Two Sigma Connect: Rental Listing Inquiries},
\href{https://www.census.gov/programs-surveys/acs/microdata.html}{Census ACS Public Use Microdata}.
\end{enumerate}
\end{frame}
